\subsection{Step 3: Process Artifacts Mining}

\begin{enumerate}

    \item \textbf{OCPM / BPMN Process Mining.}
    The per-role OCEL partitions served as input for two complementary discovery steps, both executed with the \texttt{pm4py} library.
    Object-Centric Directly-Follows Graphs (OC-DFGs) were computed for each role log, capturing activity sequences together with the multi-typed objects—issues, commits, pull requests, and tasks—co-involved in each event~\cite{notebook:05}.
    For roles exhibiting sufficient behavioural diversity, notably the \texttt{issue\_reporter} (16{,}037 events, 34 distinct activities) and \texttt{bot} (3{,}394 events, 11 activities), frequency-annotated OC-DFG visualisations were produced and persisted to \texttt{role\_logs/}.
    Complementing this, imperative BPMN models were derived via the Inductive Miner algorithm: each role log was flattened onto a primary object type, converted to a traditional event log, and the resulting process tree was transformed into a BPMN diagram exported as both a structured \texttt{.bpmn} file and a rendered image~\cite{notebook:08}.
    Together, these artefacts provide a per-role view of \emph{how} work is carried out, serving as the imperative process layer of the pipeline.

    \item \textbf{LLM Process Elicitation.}
    To produce human-readable process narratives, each role-specific OCEL file was first analysed programmatically to extract a structured process profile comprising activity frequencies, object-type distributions, directly-follows flow patterns, and event--object interaction counts~\cite{notebook:07}.
    This structured summary was then submitted to a large language model via a prompt that prescribed a fixed section schema—covering process overview, main activities, object interactions, flow patterns, and a concise business interpretation—ensuring consistent and comparable output across roles.
    The model generated a comprehensive markdown document for each of the eight roles, grounding its narrative directly in the quantitative patterns extracted from the event log rather than from general domain knowledge alone.
    The resulting per-role process descriptions were persisted as markdown files in \texttt{role\_logs/} and subsequently used to inform the role-grounding prompts of the agentic framework.

    \item \textbf{Declarative Mining.}
    Declarative process mining was applied to the role-specific logs to elicit normative behavioural constraints, complementing the imperative models produced in the previous steps~\cite{notebook:09}.
    Using the \texttt{Declare4Py} library's \texttt{DeclareMiner}, each log was first flattened onto the \texttt{issue} object type to yield a case-centric representation, and constraint discovery was executed with a minimum support threshold of 0.5 and a maximum constraint cardinality of two.
    Roles exhibiting fewer than two distinct activities—spanning \texttt{contributor}, \texttt{devops\_engineer}, \texttt{feature\_developer}, \texttt{maintainer}, \texttt{quality\_engineer}, and \texttt{technical\_writer}—were excluded from discovery due to insufficient activity diversity for meaningful constraint inference.
    Of the two eligible roles, the \texttt{issue\_reporter} log yielded three DECLARE constraints (\texttt{Existence1}, \texttt{Absence2}, and \texttt{Exactly1}) centred on the \texttt{closed} activity, encoding a normative expectation that issue closure occurs exactly once per case; the \texttt{bot} log, despite eleven distinct activities, produced no constraints above the support threshold.
    The discovered models were serialised as JSON artefacts in \texttt{declare\_models/} and subsequently integrated into the agentic framework's policy guard component as runtime behavioural guardrails.

\end{enumerate}
